% =============================================================================
%  Backward Chaining
% =============================================================================
\section{Backward Chaining}
\label{sec:backward_chaining}

\subsection{Definition}

Backward chaining is a \emph{goal-driven} inference algorithm.
It begins with a query and tries to justify it from the knowledge base.
Rather than deriving every fact that might be true, the algorithm asks a more
direct question: what rule could prove this goal, and what would that rule
require?

The contrast with forward chaining is mostly about direction.
Forward chaining starts from known facts and pushes them through rules until no
new facts appear.
Backward chaining starts from the desired conclusion and pulls it apart into
smaller goals.
A proof succeeds when those goals eventually reach facts already in the
knowledge base.
If one required sub-goal has no supporting fact or rule, that proof branch
fails.

% =============================================================================
\subsection{Horn Clause Knowledge Base}

A Horn knowledge base stores facts and rules as Horn clauses.
Each rule has one conclusion, called the \emph{head}, and zero or more
conditions, called the \emph{body}.
For backward chaining, that single head is the important part.
The prover receives a goal, looks for rules whose head matches it, and then
uses the rule body as the next set of goals.

\begin{equation}
  H \leftarrow B_1 \land B_2 \land \cdots \land B_n
\end{equation}

In Prolog-style notation:

\begin{verbatim}
head :- body1, body2, ..., bodyn.
\end{verbatim}

Logically, this means:

\begin{equation}
  B_1 \land B_2 \land \cdots \land B_n \Rightarrow H.
\end{equation}

An empty body gives a fact:

\begin{verbatim}
Parent(alice, bob).
\end{verbatim}

A non-empty body gives a rule:

\begin{verbatim}
Ancestor(x, y) :- Parent(x, y).
Ancestor(x, z) :- Parent(x, y), Ancestor(y, z).
\end{verbatim}

The first rule says that a parent is an ancestor.
The second rule handles ancestry through another person.
If the current goal is \texttt{Ancestor(alice, carol)}, the prover can match it
against the head \texttt{Ancestor(x, z)}.
After that match, the body says what remains to be shown:

\begin{quote}
To prove \texttt{Ancestor(x, z)}, prove \texttt{Parent(x, y)} and
\texttt{Ancestor(y, z)}.
\end{quote}

Horn clauses fit this proof style more directly than a general clause
knowledge base.
A CNF clause is written as a disjunction:

\begin{equation}
  L_1 \lor L_2 \lor \cdots \lor L_k.
\end{equation}

CNF is useful for forward chaining and unit propagation.
When a literal is known, clauses containing that literal are already satisfied,
and clauses containing its negation can be shortened.
If only one literal remains, the engine has found a new fact.
That representation works well when known facts are being pushed through the
whole clause set.

Backward chaining works from the other end.
It starts with a goal and asks which rule could conclude it.
A general disjunction does not give one obvious conclusion to match against;
several literals may be interpreted as candidates.
Horn clauses avoid that problem by separating the conclusion from the
conditions.
They can also be indexed by head predicate, so a prover can inspect rules for
\texttt{Ancestor(...)} without scanning unrelated rules first.

% =============================================================================
\subsection{SLD Resolution}

SLD resolution is the usual proof procedure for definite Horn clauses.
The resolver keeps a list of goals.
It takes one goal, usually the leftmost one, and searches for a clause whose
head can unify with it.
When unification succeeds, the selected goal is replaced by the clause body.
The proof then continues with this new goal list.
If the branch fails, the resolver backtracks and tries another matching
clause.

Unification is the matching step.
For example, the query

\begin{verbatim}
Ancestor(alice, z)
\end{verbatim}

unifies with the rule head

\begin{verbatim}
Ancestor(x, z)
\end{verbatim}

with the binding $x = alice$.
That binding is then applied to the rule body, and the resolver continues with
the resulting parent and ancestor sub-goals.

Before a clause is used, its variables are renamed.
This avoids accidental capture when the same rule is used more than once in a
proof.
This step is usually called \emph{standardization apart}.

% =============================================================================
\subsection{Pseudocode}

\begin{algorithm}[h]
\caption{Backward chaining over Horn clauses}
\label{alg:backward_chaining}
\begin{algorithmic}[1]
\Require Knowledge base $KB$, goal list $G$, substitution $\theta$
\Ensure A substitution proving all goals, or failure
\Procedure{BackwardChaining}{$KB, G, \theta$}
  \If{$G$ is empty}
    \State \Return $\theta$
  \EndIf
  \State $g \gets \Call{Apply}{\theta, \Call{First}{G}}$
  \State $R \gets \Call{Rest}{G}$
  \ForAll{clause $C \in KB$ whose head predicate matches $g$}
    \State $C' \gets \Call{StandardizeApart}{C}$
    \State $\theta' \gets \Call{Unify}{g, \Call{Head}{C'}, \theta}$
    \If{$\theta'$ is not failure}
      \State $G' \gets \Call{Body}{C'} \mathbin{+\!\!+} R$
      \State result $\gets \Call{BackwardChaining}{KB, G', \theta'}$
      \If{result is not failure}
        \State \Return result
      \EndIf
    \EndIf
  \EndFor
  \State \Return failure
\EndProcedure
\end{algorithmic}
\end{algorithm}

The procedure above is a depth-first search through possible proof paths.
Facts close a branch because they have no body to prove.
Rules expand the selected goal into smaller goals.
When one branch fails, the search backs up and tries another rule.

% =============================================================================
\subsection{Example}

Take a $2 \times 2$ puzzle with one clue and one inequality:

\[
\begin{array}{ccc}
1 & < & \cdot \\
\cdot &   & \cdot
\end{array}
\]

The top-left cell is already fixed to $1$.
Because it must be smaller than the top-right cell, and the only values are
$1$ and $2$, the top-right cell has to be $2$.
Backward chaining reaches that result by asking for the value directly:

\begin{verbatim}
?- Val(0, 1, v)
\end{verbatim}

The proof only needs a few Horn clauses:

\begin{verbatim}
Val(0, 0, 1).

NotVal(0, 1, 1) :- Val(0, 0, 1).
Val(0, 1, 2)    :- NotVal(0, 1, 1).
\end{verbatim}

The fact records the clue.
The first rule rules out value $1$ for the top-right cell, since $1 < 1$ is
false.
The second rule captures the tiny $2 \times 2$ domain: once value $1$ is ruled
out, value $2$ is the only value left.

The proof then runs backward from the query:

\begin{enumerate}
  \item Query \texttt{Val(0, 1, v)}.
  \item Match it with the rule head \texttt{Val(0, 1, 2)}, giving $v = 2$.
  \item Replace the original goal with \texttt{NotVal(0, 1, 1)}.
  \item Prove \texttt{NotVal(0, 1, 1)} from the clue \texttt{Val(0, 0, 1)}.
  \item Since \texttt{Val(0, 0, 1)} is already a fact, the branch succeeds.
\end{enumerate}

The returned substitution is $v = 2$.
In this small case the answer is obvious by inspection, but the proof shows
the normal backward chaining pattern: query one value, choose a rule that can
conclude it, then prove the rule body.

% =============================================================================
\subsection{Behavior}

Backward chaining works best when the program needs an answer to one specific
query.
It does not spend time deriving facts that have nothing to do with that query.
That is why the method appears so often in logic programming, expert systems,
and rule-based question answering.

The cost is repeated search.
The same sub-goal can be proved again on different branches unless the prover
stores previous results.
Recursive rules can also loop forever without some guard, such as a depth
limit, cycle detection, or tabling.

% =============================================================================
\subsection{Complexity}

Let $b$ be the average number of clauses that can match a selected goal, and
let $d$ be the maximum proof depth.
In the worst case, backward chaining explores a search tree of size

\begin{equation}
  O(b^d).
\end{equation}

The memory cost is smaller because the search is depth-first.
Ignoring the knowledge base itself, the recursion stack grows with the proof
depth:

\begin{equation}
  O(d).
\end{equation}

For definite Horn clauses, every successful proof is sound: the answer follows
from the facts and rules in the knowledge base.
Completeness depends on the search strategy.
Depth limits and other safeguards prevent runaway recursion, but they may also
cut off a valid proof.
